\documentclass[11pt,a4paper]{article}
\usepackage[margin=2.2cm]{geometry}
\usepackage{fontspec}
\setmainfont{Latin Modern Roman}
\setmonofont{Latin Modern Mono}
\usepackage[provide=*]{babel}
\babelprovide[import,main]{romanian}
\usepackage{xcolor}
\usepackage{booktabs,tabularx,array}
\usepackage{graphicx}
\usepackage{enumitem}
\usepackage{listings}
\usepackage[most]{tcolorbox}
\usepackage[hidelinks]{hyperref}
\usepackage{titlesec}
\usepackage{fancyhdr}
\usepackage{needspace}

\definecolor{accent}{HTML}{1F5F8B}
\definecolor{trap}{HTML}{B3261E}

\titleformat{\section}{\Large\bfseries\color{accent}}{}{0pt}{}[\vspace{0.2em}\color{accent}\titlerule]
\titlespacing*{\section}{0pt}{1.6em}{0.8em}
\titleformat{\subsection}{\large\bfseries\color{accent}}{}{0pt}{}
\titlespacing*{\subsection}{0pt}{1.1em}{0.4em}
\setlist{itemsep=0.25em,topsep=0.3em}
\setlength{\parindent}{0pt}
\setlength{\parskip}{0.4em}

\pagestyle{fancy}
\fancyhf{}
\renewcommand{\headrulewidth}{0.3pt}
\fancyhead[L]{Administrarea rețelelor}
\fancyhead[R]{\cursnr}
\fancyfoot[C]{\thepage}

\lstset{basicstyle=\ttfamily\small,backgroundcolor=\color{black!4},frame=single,
  rulecolor=\color{black!20},xleftmargin=0.5em,xrightmargin=0.5em,
  columns=fullflexible,keepspaces=true,showstringspaces=false,breaklines=true}

\newtcolorbox{capcana}[1][Capcane de examen]{enhanced,colback=trap!5,colframe=trap,
  fonttitle=\bfseries,title={\,#1},boxrule=0.6pt,arc=2pt,left=6pt,right=6pt}

\newcommand{\img}[3][0.85]{\IfFileExists{#2}{\begin{center}
  \includegraphics[width=#1\linewidth]{#2}\\[-0.2em]{\small\itshape #3}\end{center}}{}}
\newcommand{\cod}[1]{\texttt{#1}}
\newcolumntype{Y}{>{\raggedright\arraybackslash}X}
\newcommand{\titlu}[2]{\begin{center}{\Huge\bfseries\color{accent} #1}\\[0.4em]{\large #2}\end{center}}

% titlurile nu rămân singure la capăt de pagină
\let\oldsection\section
\renewcommand{\section}{\Needspace{9\baselineskip}\oldsection}
\let\oldsubsection\subsection
\renewcommand{\subsection}{\Needspace{6\baselineskip}\oldsubsection}

\newcommand{\cursnr}{Curs 1 -- Stiva OSI}
\begin{document}
\titlu{Stiva OSI}{Administrarea rețelelor \textbullet{} Curs 1}

\section{Modelul OSI}
\begin{itemize}
  \item Un model teoretic, cu 7 nivele independente, care comunică prin pachete (date utile $+$ date de încapsulare).
  \item \textbf{SDU} $=$ pachetul primit de nivelul $j$ de la nivelul de sus ($j+1$).
  \item \textbf{PDU} $=$ pachetul cu încapsulare trimis de nivelul $j$ către nivelul de jos ($j-1$).
\end{itemize}

\img[0.4]{Images/SDU_PDU.png}{SDU și PDU}

\begin{itemize}
  \item \textbf{Scop:} interoperabilitate între protocoale diferite. Împarte un protocol în sub-protocoale abstracte și independente.
  \item Nivelul $j$ al nodului $i$ comunică cu nivelul $j$ al nodului $k$.
  \item \textbf{Istoric (posibile întrebări de memorie):}
  \begin{itemize}
    \item Studiu lansat în \textbf{1977} de ISO împreună cu CCITT (viitorul ITU-T).
    \item Definit prima dată de \textbf{Hubert Zimmermann}, în \textbf{1978}.
    \item Arhitectura pe 7 nivele: \textbf{Charles Bachman}, Honeywell.
    \item Standardizat în \textbf{1984}: \textbf{ISO 7498} și ITU-T \textbf{X.200}.
    \item ISO/IEC 7498 are 4 părți: \textbf{-1} Basic Model, \textbf{-2} Security Architecture, \textbf{-3} Naming and Addressing, \textbf{-4} Management framework.
  \end{itemize}
\end{itemize}

\bigskip
\small
\renewcommand{\arraystretch}{1.35}
\begin{tabularx}{\linewidth}{@{}p{2.3cm} Y p{2.1cm} p{3.0cm}@{}}
\toprule
\textbf{Nivel} & \textbf{Rol din slide} & \textbf{PDU} & \textbf{Exemple}\\
\midrule
\textbf{L1 Physical} & Transmite/recepționează date \textbf{nestructurate} pe mediu (radio, electric, optic). Specifică tensiuni, viteze, distanțe, modulație, cabluri, conectori etc. & bit, simbol & SPI, I2C, 802.11, 802.15.4, 802.3\\
\textbf{L2 Data Link} & Comunicare \textbf{stabilă între 2 noduri conectate direct} (node-to-node). Poate face și flow control. & \textbf{frame} & MAC și LLC din 802.15.4, HDLC\\
\textbf{L3 Network} & Comunicare între noduri conectate \textbf{indirect} sau în \textbf{rețele diferite}; rutare, adresă unică per nod. & \textbf{packet} & IP, ICMP, IPX\\
\textbf{L4 Transport} & Legătură sigură între 2 noduri (aceeași rețea sau rețele diferite): segmentare, retransmisii, ack, multiplexare. & \textbf{segment, datagramă} & TCP, UDP (din slide 19)\\
\textbf{L5 Session} & Gestionează \textbf{sesiunile} (perioade de schimb continuu de informație) și \textbf{autentificarea}. & date structurate & RPC, PAP\\
\textbf{L6 Presentation} & Encodare, serializare, date structurate, \textbf{criptare/decriptare}, tipuri de date. & -- & BER\\
\textbf{L7 Application} & Aplicația în sine, \textbf{nu există reglementare}. & -- & http, smtp, ntp, mqtt\\
\bottomrule
\end{tabularx}
\normalsize

\subsection*{Detalii pe care le pot lua la grilă}
\begin{itemize}
  \item \textbf{L2 are 2 subnivele} (ex.\ IEEE 802):
  \begin{itemize}
    \item \textbf{MAC} $=$ politici de acces la mediu (CSMA, TDMA).
    \item \textbf{LLC} $=$ reguli de încapsulare (encodare/decodare frame-uri), detecție și corecție de erori.
    \item LLC \textbf{poate} (dar nu e obligatoriu) să ofere segmentare, frame sequencing, cereri de retransmisie.
  \end{itemize}
  \item \textbf{L3:} mesajul are minim date $+$ adresă sursă $+$ adresă destinație. Rutarea găsește o cale multi-nod. Segmentarea, detecția/corecția erorilor și retransmisia sunt \textbf{opționale}.
  \item \textbf{L4:} modelul definește \textbf{5 clase de implementare, TP0--TP4}, pentru că nu toate protocoalele pot implementa toate cerințele. \textbf{CONS} $=$ connection-oriented, \textbf{CLNS} $=$ connection-less.
\end{itemize}

\section{Stiva TCP/IP}
\begin{itemize}
  \item Se numește \textbf{Internet Protocol Suite}. Componentele principale: \textbf{TCP, UDP, IP}.
  \item \textbf{Nu respectă perfect OSI}: nu implementează toate nivelele direct și individual.
\end{itemize}

\img[0.7]{Images/OSI_TCP.png}{Stiva OSI vs.\ TCP/IP}

\subsection*{Network Access Layer $=$ L1 $+$ L2 din OSI}
\begin{itemize}
  \item Fizic diferă după mediu: cablu $=$ \textbf{Ethernet (IEEE 802.3)}, wireless $=$ \textbf{WiFi (IEEE 802.11)}.
  \item Adresare \textbf{MAC pe 6 bytes}.
  \item Oferă o interfață comună către nivelul superior (IP).
\end{itemize}

\subsection*{Internet Layer $=$ L3 complet din OSI}
\begin{itemize}
  \item Adresare prin \textbf{IP $+$ netmask}, plus \textbf{rutare}.
  \item Oferă baza pentru TCP/UDP. Mai conține \textbf{ICMP, IGMP}.
  \item \textbf{NU} oferă comunicație sigură, \textbf{NU} face retransmisii, \textbf{NU} asigură ordinea pachetelor.
\end{itemize}

\subsection*{Transport Layer}
\begin{itemize}
  \item Adresare prin \textbf{port} (număr întreg). \textbf{PORT $+$ IP $=$ SOCKET}.
\end{itemize}

\medskip
\renewcommand{\arraystretch}{1.3}
\begin{tabularx}{\linewidth}{@{}p{3.4cm} Y Y@{}}
\toprule
 & \textbf{TCP} & \textbf{UDP}\\
\midrule
Tip & connection-oriented (TP4 -- CONS) & connection-less (TP4 -- CLNS)\\
Conexiune & logică, necesară înainte de date utile, monitorizată & nu există\\
Confirmare & pentru fiecare pachet & nu\\
Retransmisie & da, pentru pachete pierdute/eronate & nu\\
Segmentare/asamblare în ordine & da & nu\\
Duplicate & le identifică și le elimină & nu\\
Ordinea & respectată & nu e asigurată\\
Rezultat & comunicație \textbf{fiabilă} & --\\
\bottomrule
\end{tabularx}

\bigskip
\begin{capcana}
\begin{itemize}
  \item \textbf{IP $\neq$ fiabil.} Fiabilitatea e la TCP, nu la Internet Layer.
  \item \textbf{TCP $=$ CONS, UDP $=$ CLNS}, ambele sub TP4.
  \item \textbf{Network Access îmbină L1 $+$ L2} (nu doar L2).
\end{itemize}
\end{capcana}
\end{document}
